\begin{abstract}

The rapid proliferation of IoT edge applications demands compact, energy-efficient SoCs capable of operating under strict constraints. To avoid the hardware overhead of conventional cryptographic ISA or pipeline extensions, this paper presents a RISP-based SoC framework integrating a minimalist single-cycle processor core with dedicated AES, SHA-3, and RSA cryptographic engines via an AXI4-Lite interface. Implemented on a Xilinx Zynq-7000 FPGA, the RISP core minimizes control logic to just $1,120\text{ }\text{LUTs}$ and $170\text{ }\text{FFs}$. Benchmarked against RV32I and Ibex baselines, the RISP-based SoC reduces total power to $0.413\text{ }\text{W}$ ($0.304\text{ }\text{W}$ dynamic) and drastically cuts cycle counts across all cryptographic workloads. Specifically, it achieves $58.18 Mb/s$ in AES-128 ECB mode ($88\text{ }\text{cycles}$), completes 16-block SHA-3-256 in $6,871\text{ }\text{cycles}$, and delivers $3314.04\text{ }\mu\text{J/bit}$ energy efficiency with $301.75$ verifications per second per Watt ($\text{verifications} \cdot \text{s}^{-1} \cdot \text{W}^{-1}$) for $\text{RSA}\text{-}2048$. The proposed architecture provides an optimal hardware offloading paradigm for energy-constrained secure edge computing.

\end{abstract}


\begin{IEEEkeywords}
RISC-V, RISP, Cryptographic Hardware, Secure Edge, RISP-Based SoC.
\end{IEEEkeywords}

\section{Introduction}
\label{sec:intro}

The rapid expansion of edge IoT applications necessitates hardware-assisted security under strict power, area, and cost constraints \cite{wang2021energy}, while conventional RISC-V instruction-set extensions and tightly coupled coprocessors introduce control complexity and hardware overhead that challenge scalability \cite{wang2021energy, pham2026sha, pham2024secure, ignatius2024power}.

To address these challenges, this paper proposes a minimalist SoC framework centered on a single-cycle Instruction-Subset Processor (RISP) coordinating AES, SHA-3, and RSA engines via an AXI4-Lite interface \cite{raissp2025}, completely decoupling system control from cryptographic computation to eliminate pipeline complexity and optimize silicon resource utilization.

This research makes two primary contributions:

\begin{enumerate}
    \item The RISP methodology is instantiated for memory-mapped cryptographic workloads by deriving an optimized instruction subset tailored to the target firmware.

    \item The framework is evaluated against a baseline SoC featuring a five-stage pipelined RV32IMB core under identical conditions.
\end{enumerate}

The remainder of this paper is organized as follows: Section \ref{sec:related} reviews literature on domain-specific processors and cryptographic integration; Section \ref{sec:system} details the proposed RISP-based SoC architecture and accelerator designs; Section \ref{sec:result} presents FPGA implementation results and performance comparisons; and Section \ref{sec:Conclusion} concludes the work and outlines future directions for edge IoT security.

\section{Related work}
\label{sec:related}

ISA-level research focuses on scalar cryptographic instructions and custom Instruction Set Extensions (ISE) to reduce program memory and deliver speedups of 1.5$\times$ to 8.6$\times$ \cite{cheng2023risc, marshall2021design, pham2026sha, pham2024secure, ignatius2024power, nicsanci2022symmetric}; however, integrating cryptographic logic directly into the processor datapath increases control complexity and can negatively impact the core's critical path and maximum frequency \cite{pham2024secure, ignatius2024power}.

While coprocessors and stand-alone accelerators (such as RoCC-integrated AES and high-throughput communication crypto engines) achieve massive performance gains \cite{gomes2022fac, pham2026sha, nannipieri2021risc, nicsanci2022symmetric, ma2023design, wang2021energy}, they often retain full-featured pipelined host processors that impose unnecessary resource overhead for extreme-edge controllers \cite{simola2024risc, cheng2023risc, ma2023design, gomes2022fac}.

Memory-mapped integration via AXI4-Lite or AXI-Stream allows modular, independent accelerator verification \cite{axinte2025embedded, simola2024risc, nasahl2021hector}, while Instruction-Subset Processor (RISP) methodologies generate application-specific single-cycle cores that reduce area by $8\%$ to $43\%$ and power by $3\%$ to $30\%$ compared to full-ISA processors \cite{raissp2025}, providing the foundation for our pairing of a minimalist RISP host with memory-mapped AES, SHA-3, and RSA accelerators.

\section{Proposed RISP-Based SoC Architecture}
\label{sec:system}

The proposed SoC pairs a minimalist RISP host with AES, SHA-3, and RSA engines via an interconnect, decoupling control from security logic for edge efficiency \cite{raissp2025}.

\subsection{System Overview and Memory-Mapped Integration}

\begin{figure}[!t]
    \centering
    \includegraphics[width=1\linewidth]{image/SoC/RISP-SOC.png}
    \caption{Overall architecture of the RISP-based SoC and internal processor microarchitecture.}
    \label{fig:risp_soc}
\end{figure}

The SoC utilizes an AXI4-Lite interconnect for low-latency communication between the RISP host and memory-mapped cryptographic accelerators \cite{gomes2022fac}, offloading intensive computations to reduce host duty cycles \cite{ma2023design}. Fig.~\ref{fig:risp_soc} illustrates the architectural framework, featuring the main CPU host coordinated via an interconnect with AES, SHA-3, and RSA accelerators, an on-chip memory hierarchy, standard I/O interfaces \cite{hoang2020low}, and a DMA controller.

\subsection{Minimalist RISP Host Core}

As illustrated in Fig.~\ref{fig:risp_soc}, the single-cycle, non-pipelined RISP host executes only the instruction subset required for firmware, minimizing pipeline overhead, TCB, and logic resources to provide a verified root-of-control \cite{raissp2025}.

\begin{table}[h!]
    \centering
    \caption{Workload-specific instructions retained or omitted in the RISSP host.}
    \label{tab:instruction_types}
    % Thiết lập bảng chiếm 50% độ rộng dòng văn bản
    \begin{tabularx}{0.5\textwidth}{l X X}
        \toprule
        \textbf{Type} & \textbf{Retained Instructions} & \textbf{Omitted Instruction}\\
        \midrule
        B-type & BEQ, BLT & BNE, BGE, BLTU, BGEU \\
        R-type & XOR & ADD, SUB, AND, OR, SLL, SRL, SRA, SLT, SLTU \\
        I-type & ADDI, ANDI & SLTI, SLTIU, XORI, ORI, SLLI, SRLI, SRAI \\
        Loads & LW & LB, LBU, LH, LHU \\
        Stores & SW & SB, SH \\
        U-Types & LUI & AUIPC \\
        Jump & JAL & JALR \\
        \bottomrule
    \end{tabularx}
\end{table}



Table~\ref{tab:instruction_types} details the workload-specific RISP instructions grouped by format, satisfying the set-completeness criterion:
\begin{equation}
\mathcal{I}_{req} \subseteq \mathcal{I}_{impl} \quad \text{and} \quad \mathcal{I}_{omit} \cap \mathcal{I}_{req} = \emptyset,
\end{equation}
where $\mathcal{I}_{req}$, $\mathcal{I}_{impl}$, and $\mathcal{I}_{omit}$ denote required, implemented, and omitted instructions, respectively. By retaining essential instructions that prioritize word-aligned access and integer arithmetic, while omitting redundant opcodes, the design significantly reduces sequential logic and control complexity compared to general-purpose baselines.

\subsection{AES Accelerator Architecture}

\begin{figure}[!t]
    \centering
    \includegraphics[width=1\linewidth]{image/crypto/AES_BD.png}
    \caption{Block diagram of the AES-128 datapath core and top-level control logic.}
    \label{fig:aes_bd}
\end{figure}

To ensure confidentiality with minimal area, the rolled AES engine employs composite field S-boxes and interfaces with the RISP host via 128-bit memory-mapped registers \cite{ignatius2024power, simola2024risc}. As depicted in Fig.~\ref{fig:aes_bd}, memory-mapped CSRs and staging registers manage inputs before mode logic routes blocks through cipher modes (ECB, CBC, CTR, CFB-128) into the core engine, concluding with status flags and ciphertext storage \cite{simola2024risc}.

Let $\mathcal{E}_K(\cdot)$ denote the 128-bit AES encryption under key $K$, where $P_i$ and $C_i$ are the $i$-th plaintext and ciphertext blocks ($i \ge 1$).
Let $IV$ be the initialization vector, $T_i$ the counter state, and $\oplus$ the bitwise XOR.
The formulas for the five supported modes are:

\begin{align}
    C_i &= \mathcal{E}_K(P_i) && \text{(ECB)} \\
    C_i &= \mathcal{E}_K(P_i \oplus C_{i-1}), \; C_0 = IV && \text{(CBC)} \\
    C_i &= P_i \oplus \mathcal{E}_K(T_i), \; T_i = T_0 + i - 1 && \text{(CTR)} \\
    C_i &= P_i \oplus \mathcal{E}_K(C_{i-1}), \; C_0 = IV && \text{(CFB)}
\end{align}

\subsection{SHA-3 Accelerator Architecture}

\begin{figure}[!t]
    \centering
    \includegraphics[width=1\linewidth]{image/crypto/SHA3-BD.png}
    \caption{High-level architecture and F-permutation processing pipeline of the SHA3 accelerator.}
    \label{fig:sha3_bd}
\end{figure}

The multimode SHA-3 engine provides integrity via a sponge construction and configurable buffers \cite{huynh2024multimode}. As illustrated in Fig.~\ref{fig:sha3_bd}, an input controller handles handshaking \cite{huynh2024multimode}, a padder formats 1088-bit blocks \cite{ma2023design}, and an internal $F_{\text{permutation}}$ engine executes 24 Keccak-f rounds over a 1600-bit state before output extraction \cite{huynh2024multimode, ma2023design}.

Let $M$ be the input message, padded into blocks $m_i$ of size $r = 1088$ bits. The internal state $S$ is initialized to a 1600-bit zero array. The sponge operation and 5-step round permutation ($\theta, \rho, \pi, \chi, \iota$) are formulated as:

\begin{align}
    \text{Padding:} \quad & S^{(0)} = 0^{1600}, \quad m_i = \text{Pad}(M) \\
    \text{Absorbing:} \quad & S^{(i)} = \mathcal{F}\left(S^{(i-1)} \oplus (m_i \parallel 0^{1600-r})\right) \\
    \text{Permutation:} \quad & \mathcal{F} = \iota \circ \chi \circ \pi \circ \rho \circ \theta \\
    \text{Squeezing:} \quad & Z_i = \text{Truncate}_d(S^{(i)})
\end{align}

\subsection{RSA-Oriented Modular Arithmetic Engine}

\begin{figure}[!t]
    \centering
    \includegraphics[width=1\linewidth]{image/crypto/RSA-BD.png}
    \caption{RSA cryptographic accelerator architecture featuring a Montgomery multiplication core.}
    \label{fig:rsa_bd}
\end{figure}

A dedicated modular arithmetic engine supports RSA and ECC primitives using a word-serial Montgomery multiplier for compact, high-precision finite field calculations \cite{wang2021energy}, decoupling long-latency modular exponentiation from the host RISP to mitigate side-channel vulnerabilities and ensure consistent performance. As illustrated in Fig.~\ref{fig:rsa_bd}, an FSM controller and input staging module coordinate parameters—such as plaintext $M$, modulus $N$, and Montgomery constants—routing them via multiplexers to execute modular exponentiation \cite{wang2021energy}.

A modular arithmetic engine accelerates RSA/ECC primitives via a word-serial Montgomery multiplier \cite{wang2021energy}. As shown in Fig.~\ref{fig:rsa_bd}, an FSM controller manages staging registers to process operands ($M$, $N$) and execute exponentiation \cite{wang2021energy}.

\begin{align}
    \text{Domain Conversion:} \quad & \bar{A} = A \cdot R \pmod N \\
    \text{Montgomery Mul:} \quad & \text{MonMul}(\bar{A}, \bar{B}) = \bar{A}\bar{B}R^{-1} \pmod N \\
    \text{Exponentiation:} \quad & C = M^d \pmod N
\end{align}

\subsection{Hardware-Software Control Flow}

\begin{figure}[!t]
    \centering
    \includegraphics[width=1\linewidth]{image/sequence/firmware_sq.png}
    \caption{Control and data exchange protocol between software and memory-mapped cryptographic hardware.}
    \label{fig:fw_sq}
\end{figure}

A lightweight register interface coordinates firmware and accelerators \cite{ma2023design}. As detailed in Fig.~\ref{fig:fw_sq}, the firmware requests operations, and the RISP host writes operands via memory-mapped transactions before asserting a CSR 'start' trigger to initiate hardware execution \cite{gomes2022fac, simola2024risc}; upon completion, a 'done' flag or interrupt notifies the host to retrieve and return results to the application layer \cite{gomes2022fac}.

\section{Experimental Evaluation and Results}
\label{sec:result}

\begin{figure*}[!t]
    \centering
    \includegraphics[width=1\linewidth]{image/evaluation/evaluated_soc.png}
    \caption{Vivado block diagram of the proposed SoC architecture.}
    \label{fig:result_soc}
\end{figure*}

This section evaluates the RISP-based SoC against RV32I and Ibex baselines on a Xilinx Zynq-7000 FPGA at 40 MHz via AXI4-Lite. Fig.~\ref{fig:result_soc} illustrates the SoC block diagram, featuring the RISP host integrated with AES, SHA-3, and RSA accelerators via an AXI interconnect.



\subsection{System-Level Resource Utilization}

\begin{table}[!t]
\centering
\caption{Hardware Resource and Performance Metrics}
\label{tab:hw_metrics_full}
\small
\resizebox{\columnwidth}{!}{%
\begin{tabular}{lccc}
\toprule
\textbf{Metric} & \textbf{RISP-Based SoC} & \textbf{RV32I-Based SoC~\cite{harris2021digital}} & \textbf{Ibex-Based SoC~\cite{gallmann2021swift}} \\
\midrule
CLB LUTs & $17,172$ & $17,632$ & $18,484$ \\
CLB FFs & $24,743$ & $26,233$ & $26,401$ \\
DSP Slices & $4$ & $4$ & $4$ \\
Core LUTs & $1,120$ & $1,584$ & $2,437$ \\
Core FFs & $170$ & $1,663$ & $1,841$ \\
Dyn. Power (W) & $0.304$ & $0.328$ & $0.336$ \\
Stat. Power (W) & $0.110$ & $0.110$ & $0.110$ \\
Total Power (W) & $0.413$ & $0.438$ & $0.446$ \\
WNS (ns) & $6.057$ & $5.249$ & $7.308$ \\
Freq. (MHz) & $52.79$ & $50.63$ & $56.50$ \\
\bottomrule
\end{tabular}%
}
\vspace{1pt}
\raggedright
\footnotesize\textit{Note: FFs = Flip-Flops; Dyn. = Dynamic; Stat. = Static; WNS = Worst Negative Slack; Freq. = Single-run slack-based frequency estimate.}
\end{table}

As summarized in Table \ref{tab:hw_metrics_full}, the proposed RISP-based SoC achieves a significantly compact and energy-efficient hardware footprint compared to standard RV32I and Ibex baselines, making it exceptionally well-suited for resource-constrained edge IoT devices. Specifically, the minimalist RISP core utilizes merely $1,120$ $LUTs$ and $170$ Flip-Flops, representing a dramatic reduction in core logic complexity relative to the RV32I ($1,584$ $LUTs$, $1,663$ $FFs$) and Ibex ($2,437$ LUTs, $1,841$ $FFs$) architectures. Furthermore, the system successfully lowers total power consumption to $0.413 W$ (with a dynamic power of just $0.304 W$), outperforming both baseline platforms while maintaining a robust operating frequency of $52.79 MHz$ and a positive worst negative slack (WNS) of $6.057 ns$. These characteristics demonstrate an optimal balance between low-power operation and hardware compactness tailored for secure edge computing environments.

\subsection{AES and SHA-3 Accelerator Evaluation}

\begin{table}[!t]
\centering
\caption{Functional Verification Vectors for AES-128 Modes}
\label{tab:aes_functional_full}
\small
\resizebox{\columnwidth}{!}{%
\begin{tabular}{l|>{\raggedright\arraybackslash}p{2.2cm}|>{\raggedright\arraybackslash}p{2.2cm}|>{\raggedright\arraybackslash}p{2.4cm}|>{\raggedright\arraybackslash}p{2.4cm}}
\toprule
\textbf{Mode} & \textbf{Input} & \textbf{IV / Counter} & \textbf{Expected Ciphertext} & \textbf{Observed Ciphertext} \\
\midrule
\textbf{ECB} & \texttt{$00112233$ \newline $44556677$ \newline $8899aabb$ \newline $ccddeeff$} & -- & \texttt{$69c4e0d8$ \newline $6a7b0430$ \newline $d8cdb780$ \newline $70b4c55a$} & \texttt{$69c4e0d8$ \newline $6a7b0430$ \newline $d8cdb780$ \newline $70b4c55a$} \\
\midrule
\textbf{CBC} & \texttt{$30c81c46$ \newline $a35ce411$ \newline $e5fbc119$ \newline $1a0a52ef$} & \texttt{$5086cb9b$ \newline $507219ee$ \newline $95db113a$ \newline $917678b2$} & \texttt{$73bed6b8$ \newline $e3c1743b$ \newline $7116e69e$ \newline $22229516$} & \texttt{$73bed6b8$ \newline $e3c1743b$ \newline $7116e69e$ \newline $22229516$} \\
\midrule
\textbf{CTR} & \texttt{$30c81c46$ \newline $a35ce411$ \newline $e5fbc119$ \newline $1a0a52ef$} & \texttt{$f0f1f2f3$ \newline $f4f5f6f7$ \newline $f8f9fafb$ \newline $fcfdff01$} & \texttt{$5ae4df3e$ \newline $dbd5d35e$ \newline $5b4f0902$ \newline $0db03eab$} & \texttt{$5ae4df3e$ \newline $dbd5d35e$ \newline $5b4f0902$ \newline $0db03eab$} \\
\midrule
\textbf{CFB} & \texttt{$f69f2445$ \newline $df4f9b17$ \newline $ad2b417b$ \newline $e66c3710$} & \texttt{$26751f67$ \newline $a3cbb140$ \newline $b1808cf1$ \newline $87a4f4df$} & \texttt{$c04b0535$ \newline $7c5d1c0e$ \newline $eac4c66f$ \newline $9ff7f2e6$} & \texttt{$c04b0535$ \newline $7c5d1c0e$ \newline $eac4c66f$ \newline $9ff7f2e6$} \\
\bottomrule
\end{tabular}%
}
\vspace{2pt}
\raggedright
\footnotesize\textit{Note: Complete 128-bit input, IV, counter, and ciphertexts for AES modes.}
\end{table}

\begin{table}[!t]
\centering
\caption{Functional Verification Vectors for SHA-3-256}
\label{tab:sha3_functional_full}
\small
\resizebox{\columnwidth}{!}{%
\begin{tabular}{c|c|>{\raggedright\arraybackslash}p{3.2cm}|>{\raggedright\arraybackslash}p{3.2cm}}
\toprule
\textbf{N blocks} & \textbf{Message (Bytes)} & \textbf{Message Summary (Hex)} & \textbf{Output Hash (Hex)} \\
\midrule
$1$ & $128$ & \texttt{$6bc1bee2$ $2e409f96$ \newline $e93d7e11$ $7393172a$ $...$} & \texttt{$d1985c30$ $73f0ff34$ \newline $c1717893$ $0e28f04e$ $...$} \\
\midrule
$4$ & $536$ & \texttt{$6bc1bee2$ $2e409f96$ \newline $6bc1bee3$ $2e409f97$ $...$} & \texttt{$7df3d9cd$ $8212808e$ \newline $b2a3ac20$ $38893a49$ $...$} \\
\midrule
$16$ & $2168$ & \texttt{$6bc1bee2$ $2e409f96$ \newline $6bc1bee3$ $2e409f97$ $...$} & \texttt{$c14e1d9a$ $4b135dc8$ \newline $b88f8272$ $0d4088bb$ $...$} \\
\bottomrule
\end{tabular}%
}
\vspace{2pt}
\raggedright
\footnotesize\textit{Note: Variable-length message blocks and output hashes.}
\end{table}

To validate the correctness of the proposed cryptographic accelerators, extensive functional verification was conducted using standard NIST test vectors across multiple operational modes and message lengths. As detailed in Fig.~\ref{tab:aes_functional_full} and Fig.~\ref{tab:sha3_functional_full}, the AES-128 engine successfully processes blocks across Electronic Codebook (ECB), Cipher Block Chaining (CBC), Counter (CTR), and Cipher Feedback (CFB) modes with matching expected and observed ciphertexts, while the SHA-3-256 accelerator efficiently handles variable-length message loads ranging from $1$ to $16$ blocks ($128$ to $2168$ bytes) to produce compliant hash outputs.

\begin{table}[!t]
\centering
\caption{AES-128 and SHA-3 Performance Comparison}
\label{tab:crypto_perf}
\small
\begin{tabular}{l l c c c}
\toprule
\textbf{Engine} & \textbf{Mode/Load} & \textbf{RISP} & \textbf{RV32I~\cite{harris2021digital}} & \textbf{Ibex~\cite{gallmann2021swift}} \\
\midrule
\multirow{2}{*}{AES} & ECB (Mb/s) & $58.18$ & $55.05$ & $56.89$ \\
                     & ECB (nJ/b) & $7.10$ & $7.95$ & $7.84$ \\
                     & non-ECB (Mb/s) & $41.29$ & $39.69$ & $40.63$ \\
                     & non-ECB (nJ/b) & $10.00$ & $11.03$ & $10.97$ \\
\midrule
\multirow{2}{*}{SHA-3} & 1-block (Mb/s) & $83.25$ & $81.92$ & $82.58$ \\
                       & 16-block (Mb/s/W) & $244.48$ & $213.5*$ & $209.7*$ \\
\bottomrule
\end{tabular}
\vspace{1pt}
\raggedright
\footnotesize\textit{Note: Mb/s = Megabits per second; nJ/b = Nanojoules per bit; Mb/s/W = Throughput per watt; ECB = Electronic Codebook mode; Non-ECB = CBC/CTR/CFB modes. *Values derived from comparative baseline analysis.}
\end{table}

Table \ref{tab:crypto_perf} presents the performance and energy efficiency evaluation of the cryptographic engines integrated into the RISP-based SoC compared to RV32I and Ibex baselines. For AES-128, the proposed architecture achieves a superior throughput of $58.18\text{ Mb/s}$ in Electronic Codebook (ECB) mode and $41.29\text{ Mb/s}$ in Non-ECB (CBC/CTR/CFB) modes, outperforming both baseline processors while maintaining lower energy consumption of $7.10\text{ nJ/b}$ (ECB) and $10.00\text{ nJ/b}$ (Non-ECB). Furthermore, for SHA-3-256 workloads, the core delivers an optimized 1-block throughput of $83.25\text{ Mb/s}$, demonstrating robust hardware acceleration capabilities tailored for secure edge computing.

\subsection{RSA Asymmetric Performance}

\begin{table}[!t]
\centering
\caption{Functional Verification Vectors for $\text{RSA}\text{-}2048$ Engine}
\label{tab:rsa_functional_full}
\small
\resizebox{\columnwidth}{!}{%
\begin{tabular}{l|>{\raggedright\arraybackslash}p{2.2cm}|c|>{\raggedright\arraybackslash}p{2.2cm}|>{\raggedright\arraybackslash}p{2.2cm}}
\toprule
\textbf{Module} & \textbf{Modulus $N$ (Hex)} & \textbf{Exp $E$} & \textbf{Input $M$ (Hex)} & \textbf{Output $C = M^E \pmod N$ (Hex)} \\
\midrule
\textbf{$\text{RSA}\text{-}2048$} &
\texttt{$e49f31a1$ $ff172739$ \newline $2d37e0d6$ $e7f0d8dc$ \newline $\dots$} &
\texttt{$65537$} &
\texttt{$4d1a8fd6$ $6d1ed78c$ \newline $3f0230f2$ $bfc14dd0$ \newline $\dots$} &
\texttt{$b3bc622e$ $5476ecf0$ \newline $0955af40$ $215a31d3$ \newline $\dots$} \\
\bottomrule
\end{tabular}%
}
\vspace{2pt}
\raggedright
\footnotesize\textit{Note: 2048-bit modulus, exponent, message, and cipher output vectors.}
\end{table}

\begin{table}[!t]
\centering
\caption{$\text{RSA}\text{-}2048$ Verification Performance}
\label{tab:rsa_perf}
\small
\begin{tabular}{l  c c c}
\toprule
\textbf{Metric} & \textbf{RISP} & \textbf{RV32I~\cite{harris2021digital}} & \textbf{Ibex~\cite{gallmann2021swift}} \\
\midrule
Cycles (Pure) & $320,972$ & $321,039$ & $320,974$ \\
Energy/bit ($\mu$J) & $3,314.04$ & $3,515.38$ & $3,578.86$ \\
Throughput (verify/s) & $124.62$ & $124.60$ & $124.62$ \\
Throughput/W & $301.75$ & $284.47$ & $279.42$ \\
\bottomrule
\end{tabular}
\end{table}

To validate the correctness and performance of the asymmetric cryptographic engine, $\text{RSA}\text{-}2048$ evaluations were conducted using standard test vectors. As illustrated in Table \ref{tab:rsa_functional_full}, the functional verification vectors demonstrate proper execution for the $2048\text{-bit}$ modulus, exponent, and cipher outputs. Furthermore, Table \ref{tab:rsa_perf} presents the performance and energy metrics compared against RV32I and Ibex baselines, showing that the RISP-Based SoC architecture achieves an optimized cycle count of $320,972$ cycles and a superior energy efficiency of $3,314.04\text{ }\mu\text{J/bit}$ (outperforming the $3,515.38\text{ }\mu\text{J/bit}$ of RV32I and $3,578.86\text{ }\mu\text{J/bit}$ of Ibex) while delivering a leading throughput-per-watt of $301.75\text{ verifications/s/W}$ at a consistent throughput of $124.62\text{ verifications per second}$.


\subsection{Cycle Count and Detailed Performance Comparison}

\begin{table}[!t]
\centering
\caption{Comprehensive Performance and Cycle Comparison}
\label{tab:crypto_summary_comparison}
\small
\resizebox{\columnwidth}{!}{%
\begin{tabular}{l l c c c}
\toprule
\textbf{Engine} & \textbf{Operation / Mode} & \textbf{RISP} & \textbf{RV32I~\cite{harris2021digital}} & \textbf{Ibex~\cite{gallmann2021swift}} \\
\midrule
\multirow{4}{*}{AES-128}
 & ECB (Cycles) & 88 & 93 & 90 \\
 & non-ECB (Cycles) & 124 & 129 & 126 \\
 & Throughput (ECB/non-ECB) & 58.18 / 41.29 & 55.1 / 39.7 & 56.9 / 40.6 \\
 & Energy (ECB/non-ECB) & 7.10 / 10.00 & 7.95 / 11.03 & 7.84 / 10.97 \\
\midrule
\multirow{3}{*}{SHA-3}
 & 1-block (Cycles) & $492$ & $500$ & $496$ \\
 & 1-block (Mb/s) & $83.25$ & $81.92$ & $82.58$ \\
 & 4-block (Cycles) & $1,771$ & $1,911$ & $1,907$ \\
 & 4-block (Mb/s) & $96.85$ & $89.75$ & $89.93$ \\
 & 16-block (Cycles) & $6,871$ & $7,419$ & $7,415$ \\
 & 16-block (Mb/s) & $100.97$ & $93.51$ & $93.56$ \\
\midrule
\multirow{2}{*}{$\text{RSA}\text{-}2048$}
 & Verify (Cycles) & $320,972$ & $321,039$ & $320,974$ \\
 & Energy / Thpt/W & $3,314$ / $302$ & $3,515$ / $284$ & $3,579$ / $279$ \\
\bottomrule
\end{tabular}%
}
\vspace{2pt}
\raggedright
\footnotesize\textit{Note:  ECB = Electronic Codebook mode; Non-ECB = CBC/CTR/CFB modes; AES cycles are per 128-bit block. SHA-3 blocks correspond to 128 B, 536 B, and 2,168 B payloads. AES throughput in Mb/s; energy in nJ/bit. RSA energy in $\mu$J/bit and throughput per watt.}
\end{table}

To provide a comprehensive evaluation of the hardware acceleration efficiency, Table~\ref{tab:crypto_summary_comparison} summarizes the execution cycles, throughput, and energy metrics across the RISP-Based SoC architecture and the RV32I and Ibex baseline SoCs for AES-128, SHA-3-256, and $\text{RSA}\text{-}2048$ workloads. For symmetric AES-128 operations, RISP-Based SoC reduces the execution overhead to $88$ cycles per block in ECB mode and $124$ cycles in non-ECB modes, outperforming both baseline processors while delivering higher throughput (up to $58.18 \text{ Mb/s}$) and lower energy consumption ($7.10 \text{ nJ/bit}$). In the hash-based SHA-3-256 workloads, the custom sponge data path maintains superior cycle efficiency across 1-block ($492$ cycles), 4-block ($1,771$ cycles), and 16-block ($6,871$ cycles) message loads—completing a 16-block payload significantly faster than the $7,419$ cycles required for RV32I and $7,415$ cycles for Ibex. Finally, for the asymmetric $\text{RSA}\text{-}2048$ verification, all architectures require comparable pure cycle counts (around $320,972$ cycles) due to the inherent complexity of modular exponentiation; however, RISP-Based SoC achieves a superior energy efficiency of $3,314 \text{ }\mu\text{J/bit}$ and a throughput-per-watt of $302$ verifications/s/W. Overall, these results demonstrate that the integrated hardware accelerators effectively minimize instruction overhead and enhance energy-performance trade-offs across diverse cryptographic algorithms.

\section{Conclusion}
\label{sec:Conclusion}

In this work, a secure and modular SoC architecture for extreme-edge IoT was developed, featuring a single-cycle RISP host coupled with AES, SHA-3, and RSA accelerators via AXI4-Lite. The decoupling of control from cryptographic operations ensured a deterministic root-of-control alongside a minimal hardware overhead. FPGA evaluations demonstrated significant resource and power savings over general-purpose alternatives, validating the effectiveness of the proposed hardware-accelerated security framework.

\bibliographystyle{IEEEtran}
\bibliography{reference}
\end{document}
